\documentclass[11pt]{article}
\usepackage[margin=1in]{geometry}
\usepackage{enumitem}
% =============================================================================
% Preambles/header.tex — shared LaTeX/Beamer preamble for this project
%
% Use from a Beamer lecture by prepending:
%     \input{header}
% and compiling with TEXINPUTS=../Preambles:$TEXINPUTS (the /compile-latex
% skill does this automatically).
%
% PALETTE CONTRACT. Color names defined here MUST match the names in
% Quarto/theme-template.scss so Beamer and Quarto renderings use the same
% palette. The scripts/check-palette-sync.sh script enforces this.
%
% To customize: edit the HEX values below AND the matching $variable in
% Quarto/theme-template.scss, then run scripts/check-palette-sync.sh.
% =============================================================================

% -----------------------------------------------------------------------------
% Engine detection + encoding
%
% Our /compile-latex skill uses XeLaTeX, where inputenc is unnecessary and
% can produce encoding warnings. Only load inputenc under pdfTeX.
% -----------------------------------------------------------------------------
\usepackage{iftex}
\ifPDFTeX
  \usepackage[utf8]{inputenc}
\fi

\usepackage{amsmath, amssymb, amsthm}
\usepackage{booktabs}
\usepackage{graphicx}

% -----------------------------------------------------------------------------
% PALETTE — mirrors Quarto/theme-template.scss variable names.
% Load xcolor and define colors BEFORE anything that references them
% (notably hyperref, hypersetup, and the Beamer color assignments).
% -----------------------------------------------------------------------------
\usepackage{xcolor}

\definecolor{primary-blue}{HTML}{012169}      % headings, accents
\definecolor{primary-gold}{HTML}{B9975B}      % emphasis, borders
\definecolor{highlight-yellow}{HTML}{F2A900}  % markers, alerts
\definecolor{light-bg}{HTML}{E8EDF5}          % title / section backgrounds
\definecolor{jet}{HTML}{1A1A1A}               % body text

% Semantic colors (used in diagrams and annotations)
\definecolor{positive}{HTML}{15803D}          % good / observed / identified
\definecolor{negative}{HTML}{B91C1C}          % bad / problematic / bias
\definecolor{neutral}{HTML}{525252}           % reference / context

% Highlight accents (match .hi-* classes in the SCSS)
\definecolor{hi-slate}{HTML}{314F4F}
\definecolor{hi-green}{HTML}{2E8B57}
\definecolor{hi-red}{HTML}{C41E3A}

% Semantic aliases so skill templates and snippets can write
% \textcolor{accent}{...} without hard-coding "primary-blue".
\colorlet{accent}{primary-blue}
\colorlet{accent2}{primary-gold}

% -----------------------------------------------------------------------------
% Hyperref — loaded AFTER xcolor + palette so link colors resolve.
% -----------------------------------------------------------------------------
\usepackage{hyperref}
\hypersetup{colorlinks=true, linkcolor=primary-blue, urlcolor=primary-blue,
            citecolor=primary-blue, pdfborder={0 0 0}}

% -----------------------------------------------------------------------------
% Course code — set once per deck (after \input{header}, before \begin{document})
% via \coursecode{CS401}. Shown in the footer on every slide so a deck's course
% is identifiable without opening the title page. Empty by default.
% -----------------------------------------------------------------------------
\makeatletter
\newcommand{\coursecode}[1]{\gdef\@coursecodetext{#1}}
\gdef\@coursecodetext{}
\makeatother

% -----------------------------------------------------------------------------
% Beamer theme assignments (only apply under Beamer).
%
% We detect Beamer via \@ifclassloaded{beamer}, which is the documented,
% non-internal way. This must be wrapped in \makeatletter/\makeatother
% because \@ifclassloaded contains an @-catcode character.
% -----------------------------------------------------------------------------
\makeatletter
\@ifclassloaded{beamer}{%
  \usefonttheme{professionalfonts}
  \setbeamercolor{structure}{fg=primary-blue}
  \setbeamercolor{frametitle}{fg=primary-blue}
  \setbeamercolor{title}{fg=primary-blue}
  \setbeamercolor{subtitle}{fg=primary-gold}
  \setbeamercolor{author}{fg=primary-blue}
  \setbeamercolor{institute}{fg=neutral}
  \setbeamercolor{date}{fg=primary-gold}
  \setbeamercolor{itemize item}{fg=primary-blue}
  \setbeamercolor{itemize subitem}{fg=primary-gold}
  \setbeamercolor{alerted text}{fg=primary-gold}
  \setbeamercolor{block title}{fg=white, bg=primary-blue}
  \setbeamercolor{block body}{bg=light-bg}
  % Semantic block variants: block=definition/concept (blue), exampleblock=
  % worked example (green/positive), alertblock=key takeaway or caution (gold).
  % Keep to <=2 colored boxes per slide across ALL three variants (INV-7).
  \setbeamercolor{block title example}{fg=white, bg=positive}
  \setbeamercolor{block body example}{bg=light-bg}
  \setbeamercolor{block title alerted}{fg=white, bg=primary-gold}
  \setbeamercolor{block body alerted}{bg=light-bg}

  % Minimal footer: course code (left) + page X / N (right), no navigation clutter
  \setbeamertemplate{navigation symbols}{}
  \setbeamertemplate{footline}{%
    \color{neutral}\footnotesize\hspace{0.7em}\@coursecodetext\hfill
    \insertframenumber/\inserttotalframenumber\hspace{0.7em}\vspace{0.3em}}
}{}
\makeatother

% -----------------------------------------------------------------------------
% TikZ setup — libraries the snippet gallery uses
% -----------------------------------------------------------------------------
\usepackage{tikz}
\usetikzlibrary{arrows.meta, positioning, calc, decorations.pathreplacing,
                fit, shapes.geometric, backgrounds}

% Shared TikZ styles — snippets and hand-written diagrams can pick these up
% via `\begin{tikzpicture}[every node/.style={...}]` or by naming specific
% styles (e.g., `\node[dag-node] (X) at (0,0) {X};`).
\tikzset{
  % Node styles for causal DAGs and similar
  dag-node/.style = {
    draw=primary-blue, fill=light-bg, rounded corners=2pt,
    minimum width=1.2cm, minimum height=0.9cm, align=center, font=\small
  },
  decision-node/.style = {
    draw=primary-gold, fill=white, diamond, aspect=1.8,
    minimum width=1.8cm, minimum height=1.2cm, align=center, font=\small,
    inner sep=1pt
  },
  % Edge styles — observed vs counterfactual
  observed-edge/.style = {-{Latex[length=2.5mm]}, thick, draw=jet},
  counterfactual-edge/.style = {-{Latex[length=2.5mm]}, thick, draw=neutral, dashed},
  confound-edge/.style = {-{Latex[length=2.5mm]}, thick, draw=negative, dashed},
  % Data-point styles for plots
  observed-dot/.style = {fill=primary-blue, draw=primary-blue, circle, inner sep=1.2pt},
  counterfactual-dot/.style = {fill=white, draw=neutral, circle, inner sep=1.2pt}
}

% -----------------------------------------------------------------------------
% Convenience macros
% -----------------------------------------------------------------------------
\newcommand{\muted}[1]{\textcolor{neutral}{#1}}
\newcommand{\key}[1]{\textcolor{primary-gold}{\textbf{#1}}}
\newcommand{\good}[1]{\textcolor{positive}{#1}}
\newcommand{\bad}[1]{\textcolor{negative}{#1}}

% Standout frame macro: full-bleed dark background for section transitions.
% Beamer-only (uses \begin{frame}/\setbeamercolor/\usebeamerfont) -- do not
% call from an article-class file (e.g. Notes/<CODE>/*-notes.tex). \muted,
% \key, \good, \bad, and \coursecode above are plain xcolor/gdef and are
% safe to use from either class; verified when Notes/article-class support
% was added.
% Expand: \transitionslide{Your transition title}
\newcommand{\transitionslide}[1]{%
  \begingroup
  \setbeamercolor{background canvas}{bg=primary-blue}
  \begin{frame}[plain]
    \centering
    \vfill
    {\usebeamerfont{title}\color{white}\Large #1\par}
    \vfill
  \end{frame}
  \endgroup
}

% Section divider frame (Beamer-only), an alternative to \transitionslide with a
% darker two-line style: near-black (jet) background, a small white label line
% above a larger gold title, and a thin gold rule along the bottom edge.
% Expand: \sectiondivider{Section 2}{Pointers}
\newcommand{\sectiondivider}[2]{%
  \begingroup
  \setbeamercolor{background canvas}{bg=jet}
  \begin{frame}[plain]
    \vspace*{3.0cm}
    {\color{white}\ttfamily\large #1\par}
    \vspace{0.5cm}
    {\color{primary-gold}\LARGE #2\par}
    \begin{tikzpicture}[remember picture, overlay]
      \draw[primary-gold, line width=2.5pt]
        ($(current page.south west)+(0,0.1)$) -- ($(current page.south east)+(0,0.1)$);
    \end{tikzpicture}
  \end{frame}
  \endgroup
}

% -----------------------------------------------------------------------------
% Channel watermark — "Concepts That Click" (YouTube).
%
% Article-class files (CompetitiveExam Q/A sets, Notes, Assignments) get a
% light diagonal draft watermark on every page plus a centered footer naming
% the channel. Beamer decks get a subtle TikZ background watermark instead
% (the footline already carries the course code). Centralized here so every
% MCQ PDF is branded without per-file edits.
% -----------------------------------------------------------------------------
\makeatletter
\@ifclassloaded{article}{%
  \usepackage{draftwatermark}
  \SetWatermarkText{Concepts That Click}
  \SetWatermarkScale{1}
  \SetWatermarkFontSize{2cm}
  \SetWatermarkLightness{0.86}
  \SetWatermarkAngle{45}
  \usepackage{fancyhdr}
  \pagestyle{fancy}
  \fancyhf{}
  \fancyfoot[C]{\footnotesize\textcolor{neutral}{Concepts That Click~|~YouTube: \href{https://www.youtube.com/@concepts.thatclick}{@concepts.thatclick}}}
  \renewcommand{\headrulewidth}{0pt}
  \renewcommand{\footrulewidth}{0pt}
  \fancypagestyle{plain}{%
    \fancyhf{}%
    \fancyfoot[C]{\footnotesize\textcolor{neutral}{Concepts That Click~|~YouTube: \href{https://www.youtube.com/@concepts.thatclick}{@concepts.thatclick}}}%
    \renewcommand{\headrulewidth}{0pt}%
    \renewcommand{\footrulewidth}{0pt}%
  }
}{}
\@ifclassloaded{beamer}{%
  \setbeamertemplate{background}{%
    \begin{tikzpicture}[remember picture, overlay]
      \node[rotate=45, opacity=0.055, align=center]
        at (current page.center)
        {\Huge\textcolor{neutral}{Concepts That Click}};
    \end{tikzpicture}%
  }
}{}
\makeatother 
\coursecode{JKSSB}
\renewcommand{\thesection}{29.\arabic{section}}
\newtheorem{example}{Example}[section]
\newenvironment{definitionbox}[1]{%
  \par\noindent\textbf{Definition (#1).}\ \itshape
}{\par\vspace{0.3em}}
\title{Access Keys and Relationships --- Detailed Lecture Notes}
\author{JKSSB Computer Awareness}
\date{\today}

\begin{document}
\maketitle

\section{Why keys and relationships matter}
A database keeps data in \textbf{tables}. A table is a grid: each
\textbf{record} is one row, and each \textbf{field} is one column that holds a
single kind of value. When a table holds many records, values such as a name,
a city, or a date can repeat, so a row cannot be recognised by those values
alone. A database therefore gives each row a reliable \textbf{identifier},
called a key, and it links rows that belong together using
\textbf{relationships}. Keys provide identity; relationships join tables that
share data; referential integrity keeps those joins correct.

\begin{definitionbox}{Table, record, field}
A \textbf{table} stores data about one subject as rows and columns. A
\textbf{record} is one complete row. A \textbf{field} is one column, holding
one attribute for every record.
\end{definitionbox}

\begin{center}
\begin{tabular}{p{0.28\textwidth}p{0.60\textwidth}}
\toprule
\textbf{Term} & \textbf{Meaning in this lesson} \\
\midrule
Primary key (PK) & A field or set of fields that uniquely identifies each record. \\
Foreign key (FK) & A field that refers to the primary key of another table. \\
Composite key & A primary key made of two or more fields together. \\
Candidate key & Any field that could serve as the primary key. \\
Alternate key & A candidate key that was not chosen as the primary key. \\
Junction table & A table that links two tables in a many-to-many relationship. \\
Referential integrity & The rule that a foreign-key value must match an existing primary-key value. \\
\bottomrule
\end{tabular}
\end{center}

The rest of the lesson uses these terms in one consistent sense. A
\textbf{primary key} identifies a record inside its own table. A
\textbf{foreign key} points from one table to another table's primary key. A
\textbf{composite key} is simply a primary key built from more than one field.
A \textbf{junction table} is the third table that implements a many-to-many
relationship.

\section{The primary key}
A \textbf{primary key (PK)} is a field, or a set of fields, that
\textbf{uniquely identifies} each record in a table. It has three firm rules.
First, its value must be \textbf{unique}: no two records may share the same
primary-key value. Second, it \textbf{cannot be left blank} (it cannot contain
a null value), because a blank value cannot identify anything. Third, each
table has \textbf{at most one} primary key, although that one key may be built
from several fields.

\begin{definitionbox}{Primary key}
A field (or set of fields) whose value is unique for every record and is never
blank, so it can identify each record in the table.
\end{definitionbox}

\begin{example}
In an Employee table, the field \texttt{EmployeeID} is a good primary key: every
employee receives a different number, and no employee is stored without one.
The employee's name would be a poor choice, because two employees can share a
name.
\end{example}

Access can generate a primary key automatically with the \textbf{AutoNumber}
data type, which inserts a unique number for each new record. A primary key is
what makes a table the \textbf{``one'' side} of a one-to-many relationship.

\section{Candidate keys, alternate keys and composite keys}
Any field that could serve as the primary key is called a \textbf{candidate
key}. When a table has more than one candidate key, only one is selected as
the primary key; the remaining candidates are called \textbf{alternate keys}.
For example, a Student table might have both \texttt{RollNumber} and
\texttt{Email} as candidate keys; if \texttt{RollNumber} is chosen, then
\texttt{Email} becomes an alternate key.

A \textbf{composite key} (also called a \textbf{compound key}) is a primary key
made of \textbf{two or more fields used together}. The combination is unique
even when each field on its own is not. Composite keys appear most often in
junction tables, where two foreign keys together form the primary key.

\begin{definitionbox}{Composite key}
A primary key that consists of two or more fields, whose combined value is
unique for every record.
\end{definitionbox}

\begin{example}
An OrderDetails table records each product line within an order. Neither
\texttt{OrderID} nor \texttt{ProductID} is unique by itself, but the pair
(\texttt{OrderID}~+~\texttt{ProductID}) is unique for every line item, so the
pair can be the primary key. This is a composite key.
\end{example}

A good primary key is also \textbf{stable} and \textbf{meaningless}: it should
not need to change later. This is why Access often uses an AutoNumber field
rather than a value such as a name or a phone number, which a person may edit.

\begin{center}
\begin{tabular}{p{0.22\textwidth}p{0.62\textwidth}}
\toprule
\textbf{Key} & \textbf{What it does} \\
\midrule
Primary key & Uniquely identifies a record; unique and not blank; at most one per table. \\
Candidate key & Any field eligible to become the primary key. \\
Alternate key & A candidate key that was not chosen as the primary key. \\
Composite key & A primary key made of two or more fields together. \\
Foreign key & A field that refers to another table's primary key; may repeat. \\
\bottomrule
\end{tabular}
\end{center}

\section{The foreign key}
A \textbf{foreign key (FK)} is a field in one table that refers to the
\textbf{primary key} of another table. The foreign key is the link that
creates a relationship between the two tables. Unlike a primary key, a foreign
key \textbf{may repeat} and \textbf{may be blank}; it is not unique. The table
that holds the foreign key is the \textbf{child} (or related) table, and the
table whose primary key is referenced is the \textbf{parent} table.

\begin{definitionbox}{Foreign key}
A field in one table whose values must match the primary-key values of another
table, used to link the two tables.
\end{definitionbox}

\begin{example}
In an Employee table, the field \texttt{DepartmentID} is a foreign key that
refers to the \texttt{DepartmentID} primary key in the Department table. Many
employees can share the same \texttt{DepartmentID}, so the foreign key repeats;
that is expected and correct.
\end{example}

\begin{center}
\begin{tabular}{p{0.24\textwidth}p{0.32\textwidth}p{0.32\textwidth}}
\toprule
\textbf{Point} & \textbf{Primary key} & \textbf{Foreign key} \\
\midrule
Purpose & Uniquely identifies a record & Links to another table \\
Uniqueness & Must be unique & May repeat \\
Blank allowed & No & Yes \\
Count per table & At most one & Can be more than one \\
\bottomrule
\end{tabular}
\end{center}

\section{Types of relationships}
A \textbf{relationship} works between two tables that share a common field ---
usually a primary key in one table and a foreign key in the other. Access
supports three relationship types: \textbf{one-to-one},
\textbf{one-to-many}, and \textbf{many-to-many}. Relationships are created and
edited in the \textbf{Relationships} window.

\begin{center}
\begin{tabular}{p{0.20\textwidth}p{0.38\textwidth}p{0.32\textwidth}}
\toprule
\textbf{Type} & \textbf{Meaning} & \textbf{How it is built} \\
\midrule
One-to-one & Each record in one table matches at most one in the other & A shared field that is unique in both tables \\
One-to-many & One record matches many records in the other table & Primary key on the ``one'' side, foreign key on the ``many'' side (most common) \\
Many-to-many & Many records match many records & Two one-to-many relationships through a junction table \\
\bottomrule
\end{tabular}
\end{center}

The rule to hold on to is simple: the \textbf{``one'' side} always contributes
the \textbf{primary key}, and the \textbf{``many'' side} always carries the
\textbf{foreign key}. A one-to-one relationship is just a special case in
which both sides are unique.

\subsection{One-to-one}
In a \textbf{one-to-one} relationship, each record in the first table matches
at most one record in the second table, and vice versa. It is used to split a
table with many fields into smaller parts, or to keep sensitive data in a
separate table that can be secured on its own. It is the \textbf{least common}
of the three types.

\begin{example}
An Employee table and an EmployeeConfidential table are linked one-to-one:
each employee has exactly one confidential record, and each confidential
record belongs to exactly one employee.
\end{example}

\subsection{One-to-many}
In a \textbf{one-to-many} relationship, one record in the first table matches
\textbf{many} records in the second table. The ``one'' side holds the
\textbf{primary key}, and the ``many'' side holds the \textbf{foreign key}.
This is the \textbf{most common} relationship type used in Access.

\begin{example}
One Department has many Employees, but each Employee belongs to only one
Department. The Department table is the ``one'' side and the Employee table is
the ``many'' side, linked by \texttt{DepartmentID}.
\end{example}

\subsection{Many-to-many and the junction table}
In a \textbf{many-to-many} relationship, many records in the first table match
many records in the second. Access \textbf{cannot} store a many-to-many
relationship directly. Instead it is split into \textbf{two one-to-many}
relationships by adding a third table called a \textbf{junction table} (also
known as an \textbf{associative} or \textbf{linking} table). The junction
table contains a foreign key for each of the two other tables.

\begin{definitionbox}{Junction table}
A table that connects two tables in a many-to-many relationship, holding a
foreign key for each of them.
\end{definitionbox}

A junction table's own primary key is usually a \textbf{composite key} formed
from the two foreign keys. It may also carry extra fields that describe the
link itself, such as a quantity, a grade, or a date.

\begin{example}
Students and Courses have a many-to-many relationship: a student takes many
courses, and a course has many students. Access implements this with an
Enrollments junction table that holds \texttt{StudentID} and \texttt{CourseID}
as its composite primary key, together with a \texttt{Grade} field.
\end{example}

\section{Creating a relationship in Access}
In the Access interface, relationships are built and edited in a small number
of steps. On the \textbf{Database Tools} tab, the \textbf{Relationships} button
opens the \textbf{Relationships} window. The tables to be linked are added to
this window, and then the \textbf{primary-key field} of one table is
\textbf{dragged onto} the matching foreign-key field of the other table.

\begin{enumerate}
  \item Open the \textbf{Relationships} window from the \textbf{Database Tools}
    tab.
  \item Add the two tables that will be linked.
  \item Drag the primary key of the ``one'' table onto the foreign key of the
    ``many'' table.
  \item In the \textbf{Edit Relationships} dialog box, choose
    \textbf{Enforce Referential Integrity}.
  \item Optionally tick \textbf{Cascade Update} and \textbf{Cascade Delete},
    then confirm.
\end{enumerate}

\begin{example}
To link Departments and Employees, drag \texttt{DepartmentID} from the
Department (parent) table onto \texttt{DepartmentID} in the Employee (child)
table, then enforce referential integrity. Access now refuses to store an
employee whose \texttt{DepartmentID} does not exist in the Department table.
\end{example}

\section{Referential integrity}
\textbf{Referential integrity} is the rule that a foreign-key value must
\textbf{match an existing primary-key value} in the related table (or be
blank). It prevents \textbf{orphan records}, which are child rows that point to
a parent row that does not exist. In Access, referential integrity is switched
on in the \textbf{Edit Relationships} dialog box (opened from the
\textbf{Relationships} window).

\begin{definitionbox}{Referential integrity}
A rule that keeps the foreign key in a child table consistent with the primary
key in the parent table: a child value must match an existing parent value or
be null.
\end{definitionbox}

When referential integrity is enforced, Access offers two related options.
\textbf{Cascade Update} means that if the parent primary-key value changes,
Access automatically updates the matching foreign-key values in the child
table. \textbf{Cascade Delete} means that if a parent record is deleted, Access
automatically deletes the matching child records as well.

\begin{example}
If a Department record is deleted and Cascade Delete is on, all Employees that
belong to that Department are deleted too. Without referential integrity,
Access would refuse to delete the Department while those Employees still point
to it.
\end{example}

\section{Access file formats: \texttt{.mdb} vs \texttt{.accdb}}
The file \textbf{extension} shows which Access format a database uses.
\texttt{.mdb} (\textit{Microsoft Database}) is the \textbf{older} Access file
format, used by Access 2003 and earlier versions. \texttt{.accdb}
(\textit{Access Database}) is the \textbf{newer, default} format introduced
with Access 2007 and used since. Access 2007 and later can still open an old
\texttt{.mdb} file, but new databases are saved as \texttt{.accdb} by default.

\begin{center}
\begin{tabular}{p{0.18\textwidth}p{0.30\textwidth}p{0.40\textwidth}}
\toprule
\textbf{Extension} & \textbf{Access version} & \textbf{Notes} \\
\midrule
\texttt{.mdb} & Access 2003 and earlier & The older format; supports user-level security. \\
\texttt{.accdb} & Access 2007 and later & The newer default; supports attachment, multi-valued and calculated fields. \\
\bottomrule
\end{tabular}
\end{center}

The newer \texttt{.accdb} format added \textbf{attachment} fields,
\textbf{multi-valued} fields, and \textbf{calculated} fields, and it integrates
better with SharePoint and Outlook. Keep the pairing straight: \texttt{.mdb}
is the \emph{old} format and \texttt{.accdb} is the \emph{new} format.
Reversing ``old'' and ``new'' is a common exam trap.

\section{Exam retrieval and common confusions}
Six distinctions clear up most errors on this topic.
\begin{enumerate}
  \item A primary key is \textbf{unique} and \textbf{not blank}; a foreign key
    \textbf{may repeat} and \textbf{may be blank}.
  \item A table has \textbf{at most one} primary key, but it can hold
    \textbf{several} foreign keys.
  \item A foreign key refers to the \textbf{primary key of another table}, not
    to its own table's primary key.
  \item \textbf{One-to-many} is the most common relationship; a
    \textbf{many-to-many} relationship is never stored directly but is built
    with a \textbf{junction table}.
  \item A junction table commonly has a \textbf{composite primary key} made
    from the two foreign keys.
  \item \texttt{.mdb} is the \textbf{older} Access format; \texttt{.accdb} is
    the \textbf{newer} format.
\end{enumerate}

\begin{example}
An item asks which key may contain duplicate values. Walk through the rules: a
primary key must be unique and not blank, so it cannot repeat; a foreign key
may repeat and may be blank; a composite key is still a primary key and so is
unique as a combination. The answer is the foreign key.
\end{example}

\begin{example}
An item states that ``a many-to-many relationship is created by adding a
foreign key to one of the two tables.'' This is wrong. A many-to-many
relationship cannot be stored in either table directly; a \textbf{junction
table} is added, and two one-to-many relationships are built through it. The
wrong option has skipped the junction table.
\end{example}

\subsection*{Quick self-check}
Before the exercise set, confirm you can state each of these in one line.
\begin{itemize}
  \item A primary key is unique, never blank, and there is at most one per
    table.
  \item A foreign key links to another table's primary key and may repeat.
  \item A composite key is a primary key of two or more fields.
  \item One-to-many is the most common relationship; the ``one'' side supplies
    the primary key and the ``many'' side the foreign key.
  \item A many-to-many relationship is always built with a junction table.
  \item Referential integrity keeps foreign-key values matching primary-key
    values and offers Cascade Update and Cascade Delete.
  \item \texttt{.mdb} is the older format; \texttt{.accdb} is the newer one.
\end{itemize}

\subsection*{Exercises}
\begin{enumerate}[label=\arabic*.]
  \item Define a primary key and state its two main rules about uniqueness and
    blank values.
  \item Explain the difference between a primary key and a foreign key.
  \item What is a composite key? Give one example where it is used.
  \item Name the three relationship types in Access and say which is the most
    common.
  \item Why does a many-to-many relationship need a junction table?
  \item State the rule of referential integrity and name its two cascade
    options.
  \item Distinguish the Access file formats \texttt{.mdb} and \texttt{.accdb}.
\end{enumerate}

\subsection*{Solutions}
\begingroup\small
\begin{enumerate}[label=\arabic*.]
  \item A primary key is a field (or set of fields) that uniquely identifies
    each record. Its value must be unique --- no two records may share it ---
    and it cannot be left blank (no null value). A table has at most one
    primary key.
  \item A primary key uniquely identifies records inside its own table and
    cannot repeat or be blank. A foreign key is a field that refers to another
    table's primary key; it may repeat and may be blank, and it creates the
    relationship between the two tables.
  \item A composite key is a primary key made of two or more fields used
    together, where the combination is unique. It is used in a junction table
    --- for example, an Enrollments table whose primary key is
    (\texttt{StudentID}~+~\texttt{CourseID}).
  \item The three types are one-to-one, one-to-many, and many-to-many. The
    \textbf{one-to-many} relationship is the most common.
  \item Access cannot store a many-to-many relationship directly. It is split
    into two one-to-many relationships by a junction table that holds a
    foreign key for each of the two tables.
  \item Referential integrity means a foreign-key value must match an existing
    primary-key value in the related table (or be blank), which prevents
    orphan records. Its two options are \textbf{Cascade Update} and
    \textbf{Cascade Delete}.
  \item \texttt{.mdb} is the older Access format used by Access 2003 and
    earlier. \texttt{.accdb} is the newer, default format used from Access
    2007 onward, adding attachment, multi-valued and calculated fields.
\end{enumerate}
\endgroup
\end{document}
